% Verbatim copy of /Users/alingalatan/workspaces/grothe/handoff/manuscript/legacy/ch25-parts/91-gaussian-point-and-cone-reach.tex (file dated 2026-10-25). Only this line was added.
% IDE module, 22 October 2026: the Gaussian point of the k-row functional, the sector Hessian, the nu-sector reduction,
% independent float checks of the contact-cone theorems, and the reach of the support-cone method.
\subsection{The Gaussian point, the sector Hessian, and the reach of support cones}
\label{sec:gaussian-point-cone-reach}

\paragraph{Status and scope.}
This section records the IDE's audit of the 22 October contact-cone
theorems (\S\ref{sec:mixed-contact-concentration},
\S\ref{sec:whole-slice-contact}) and three structural facts about the
remote functional of Conjecture~K. Formulas and reductions are
\st{P}; every number is a float quadrature \st{F} unless it is quoted
from a directed certificate. Nothing here changes the complete endpoint
or the Lean development; the Lean chain does not formalise the
original-response programme, and none of the statements below is on
its critical path (which is the sixth-moment tensor bridge, the row
reductions and the transfer link).

Throughout, $\epsilon=1891/2500$, $p=1/\epsilon$,
$L=Z(\ell+\beta(T-2))$ with $T=|Z|^2$, $\ell=1057/1250$,
$\beta=179/1250$ (the peak profile), and for a remote cubic $G$
\[
 \Lambda(G)=\E|L+G|+\E|L-G|-\frac{\|G\|_2^2}{\epsilon},
 \qquad
 F(v)=2\E|L+\sqrt v\,U|-\frac v\epsilon,\qquad S=\sup_{v\ge0}F(v),
\]
$U$ a standard complex Gaussian independent of $Z$. The circular-mean
potential of a circular law is
$\Psi_q(x)=\E|x+q|=\E\,\mathrm{M}(x,|q|)$, with
$\mathrm M(a,r)=\E_\theta|a+re^{i\theta}|=\frac2\pi(a+r)\,
\mathrm{E}\!\left(\frac{4ar}{(a+r)^2}\right)$ (complete elliptic
integral of the second kind).

\subsubsection{The Gaussian family and its closed forms}
With $x=\lambda^2/v$ and $\lambda=\sqrt T\,(\ell+\beta(T-2))$, the Rice
formulas give the one-dimensional integrals
\begin{equation}\label{eq:rice-closed-forms}
 \E\big[\,|L+\sqrt vU|\;\big|\;T\big]
 =\frac{\sqrt{\pi v}}2\,e^{-x/2}\big[(1+x)I_0(\tfrac x2)+xI_1(\tfrac x2)\big],
 \qquad
 \E\Big[\frac1{|L+\sqrt vU|}\;\Big|\;T\Big]
 =\sqrt{\frac\pi v}\,e^{-x/2}I_0(\tfrac x2),
\end{equation}
and the heat equation $\partial_v=\frac14\Delta$ with $\Delta|y|=1/|y|$
gives $F'(v)=\frac12\E|L+\sqrt vU|^{-1}-1/\epsilon$ \st{P}.
Integrating \eqref{eq:rice-closed-forms} against $e^{-T}$
(Gauss--Legendre panels, 32 nodes per panel) reproduces the rigorous
enclosure of the cubic-chaos package,
$1.3739535256614<S<1.3739535256616$, as $S=1.37395352566$ at
$v^*=0.0403448$, with $F''(v^*)=-2.26$ \st{F}. The package proves
$-4<F''<0$ on $[0,\infty)$ \st{P} (strict concavity through the heat
equation and radially decreasing densities), and the sharpness
$F(v)-6v^2/m\le\Lambda(G_{m,v})\le F(v)$ for the disjoint-product
construction $G_{m,v}=\sqrt{v/m}\sum_{j\le m}z_{3j-1}z_{3j}\bar z_{3j+1}$
\st{P}; so $S$ is the supremum over all dimensions if Conjecture~K
holds, and the proposed relaxation constant $1.3\cdot10^{-4}$ cannot be
lowered.

\subsubsection{Second variation at the Gaussian point}
\begin{lemma}[\st{P}: second variation of the modulus]
\label{lem:second-variation-modulus}
For $z\ne0$ and $w\in\C$,
$\frac{d^2}{ds^2}|z+sw|\big|_{s=0}=\Im(we^{-i\arg z})^2/|z|$.
Consequently, for $G_0$ with $|L\pm G_0|>0$ a.s.\ and a direction
$G'$ with $\E|G'|^2=1$ for which differentiation under the expectation
is justified,
\[
 \frac{d^2}{ds^2}\Lambda(G_0+sG')\Big|_{s=0}
 =\E\frac{\Im(e^{-i\theta_+}G')^2}{|L+G_0|}
 +\E\frac{\Im(e^{-i\theta_-}G')^2}{|L-G_0|}-\frac2\epsilon,
 \qquad\theta_\pm=\arg(L\pm G_0).
\]
\end{lemma}
\begin{proof}
$|z+sw|=\sqrt{(|z|+s\Re(we^{-i\arg z}))^2+s^2\Im(we^{-i\arg z})^2}$;
expand to second order. The second statement adds the two signs and
differentiates $-\|G_0+sG'\|^2/\epsilon$.
\end{proof}

Take the Gaussian point $G_0=\sigma^*W$, $W$ a diffuse circular
Gaussian remote cubic (the $N\to\infty$ limit of $G_{N,v}$), with
$v^*=\sigma^{*2}$ the maximiser of $F$. Factor the phase of $Z$ so that
$L=\lambda(T)e^{i\arg Z}$ and $W\mapsto W'$ is again standard. In the
diffuse limit the Hessian is block diagonal on the six sector directions
of the decomposition of the cubic-chaos package
($G=q(w)+(T-1)u(w)+\frac{Z^2}{\sqrt2}v(w)+ZA(w)+\bar ZB(w)+cH(Z)$,
$H(Z)=Z(T-2)/\sqrt2$): cross terms between components of different
$Z$-frequency vanish after the $Z$-phase average, except that the
A-sector $(T-1)u$ (frequency $-1$ relative to $L$) and the B-sector
$Z^2v/\sqrt2$ (frequency $+1$) are coupled through
$\E_\theta\Im(\alpha e^{-i\theta})\Im(\beta e^{i\theta})=-\frac12\Re(\alpha\beta)$,
i.e.\ through $\Re\E[uv]$, which need not vanish ($u\in\Hc10(w)$,
$v\in\Hc01(w)$ may be conjugates). The table below lists the diagonal
entries; the off-diagonal A--B entry is bounded in
\S\ref{sec:local-stability-reported} and the $2\times2$ block stays
negative definite. Writing
$w_1=\frac12\big(|\lambda+\sigma^*W'|^{-1}+|\lambda-\sigma^*W'|^{-1}\big)$
and $\psi_\pm=\arg(\lambda\pm\sigma^*W')$, the unit-norm sector
directions have \st{P}
\[
\begin{array}{ll}
 \text{remote cubic independent of }W:& h_q=\E w_1-2/\epsilon=2F'(v^*)=0,\\
 \text{the radial direction }W:&
   h_W=\E\big[\Im(e^{-i\theta_+}W')^2/|\lambda+\sigma^*W'|
            +\Im(e^{-i\theta_-}W')^2/|\lambda-\sigma^*W'|\big]-2/\epsilon,\\
 \text{A-sector }(T-1)X:& h_A=\E[(T-1)^2w_1]-2/\epsilon,\\
 \text{B-sector }Z^2\bar X/\sqrt2:& h_B=\E[T^2w_1/2]-2/\epsilon,\\
 \text{Y- and }(2,0)\text{-sectors, circular:}& h_Y=\E[Tw_1]-2/\epsilon,\\
 \text{Y-sector, real / imaginary }Y:&
   h_{Y,r}=\E\big[T(\sin^2\psi_+/|\lambda+\sigma^*W'|+\sin^2\psi_-/|\lambda-\sigma^*W'|)\big]-2/\epsilon,
   \quad h_{Y,i}=\text{(same with }\cos^2),\\
 \text{visible }\nu H(Z):&
   h_\nu=\E\big[\tfrac12T(T-2)^2(\sin^2\psi_+/|\lambda+\sigma^*W'|+\sin^2\psi_-/|\lambda-\sigma^*W'|)\big]-2/\epsilon .
\end{array}
\]
The circular ones are one-dimensional integrals by
\eqref{eq:rice-closed-forms}; the others are three-dimensional
$(T,|W'|,\arg W')$ quadratures. Their values \st{F}:
\[
\begin{array}{l|c|c}
 \text{direction}&\epsilon=1891/2500,\ (\ell,\beta)=(.8456,.1432)
 &\epsilon=149/200,\ (\ell,\beta)=(.852239,.144728)\\\hline
 q\ (\text{Gaussian family, flat})&0&0\\
 \text{radial }W&-0.366&-0.279\\
 \text{A-sector}&-0.931&-0.913\\
 \text{B-sector}&-1.867&-1.917\\
 \text{Y, }(2,0)\text{ circular}&-1.401&-1.461\\
 \text{Y real / imaginary}&-2.41\ /\ -0.395&-2.51\ /\ -0.414\\
 \nu H(Z)&-2.34&-2.45
\end{array}
\]
So the Gaussian point is a strict local maximum of $\Lambda$ in every
sector direction, the only flat direction being the Gaussian family
itself, along which $F''<0$. The smallest margins are the radial
direction (the law of $|W|$ leaving the Rayleigh law) and the imaginary
Y-sector. For comparison, at $G=0$ the exact second-order coefficients
of the same sectors are $k_j/2-p$ with the $k_j$ of the cubic-chaos
package ($A{:}-0.61$, $B{:}-1.84$, $Y{:}-1.42$, $\nu{:}-1.59$): the
structured sectors become \emph{more} negative as the remote variance
grows, because Gaussian smoothing of $1/|L+q|$ removes the mass near
$T=0$ where $1/|L|$ is large. The dangerous regime for structured
sectors is therefore small remote energy, which is where the
$G=0$ bounds of the package apply. (The package's pointwise inequality
$|l+g|+|l-g|\le2|l|+|g|^2/|l|$ is the worst orientation; for a
component whose conditional law given $Z$ is circular, the exact
second-order coefficient is half of $k_j$, which is why its
coefficient $41/20$ for the A-sector exceeds $p$ while the A-sector is
nevertheless dominated to second order at $G=0$.) This is a local
statement; it says nothing about the global maximum.

\subsubsection{The visible sector is a profile perturbation of the remote row}
\begin{lemma}[\st{P}: $\nu$-sector reduction]
\label{lem:nu-sector-reduction}
Let $q$ be a remote cubic independent of $Z$ whose law is circular, and
$\nu\in\C$. Then
\[
 \Lambda(q+\nu H)=\sum_{\pm}\E_T\,
 \Psi_q\!\Big(\sqrt T\,\Big|\ell+\beta(T-2)\pm\nu\frac{T-2}{\sqrt2}\Big|\Big)
 -\frac{\|q\|^2+|\nu|^2}{\epsilon}.
\]
\end{lemma}
\begin{proof}
$L\pm\nu H=Z\,(\ell+\beta(T-2)\pm\nu(T-2)/\sqrt2)$ has the phase of $Z$
up to the phase of the bracket, which is a function of $T$. Condition on
$T$: the phase of $Z$ is uniform and independent of $q$, and $q$ is
circular, so $\E[\,|L\pm\nu H+q|\;|\;T]=\Psi_q(\sqrt T\,|\ell+\beta(T-2)\pm\nu(T-2)/\sqrt2|)$.
Orthogonality gives $\|q+\nu H\|^2=\|q\|^2+|\nu|^2$.
\end{proof}
Hence the visible sector is exactly a perturbation of the row profile
of the remote-only functional, the object the Lean rows bound through
the gate variable, and to second order it costs
$|\nu|^2\big(\E_T[\tfrac12T(T-2)^2\,\Psi_q''(\sqrt T f(T))]-1/\epsilon\big)$
with $f=\ell+\beta(T-2)$. For the Gaussian law the bracket is
$h_\nu/2=-1.17$ and at $q=0$ it is $k_c-p=-0.29$ \st{F}. For a general
law $\Psi_q''(x)\le\E|x+q|^{-1}$, which is finite only under a small-ball
estimate for $q$ (a law with an atom at $-x_0$ makes
$\E_T\,w(T)/|x(T)-x_0|$ diverge logarithmically). This is the same input
as the two-plane small-ball estimate of the $t=0$ row and explains why
the gate price is paid in the certificates rather than derived from
moments.

\subsubsection{Independent checks of the contact-cone integrals and the slice}
The two support integrals of
Theorem~\ref{thm:mixed-contact-concentration} were recomputed by an
independent float quadrature
($I(1,0,1/4)=0.9352013$, $I(1,37/200,87/500)=0.8365651$, stable under
grid doubling) and lie $1\cdot10^{-5}$ below the directed rational caps
$467607/500000$ and $33463/40000$, as they must. With the same rows and
$u_0=47/200$, the four-endpoint envelope
\eqref{eq:mixed-contact-canonical-cap} evaluated along the slice test
$\psi_\rho=\rho X(|X|^2-2)+\sqrt{1-\rho^2}X^2\bar Y$ of
\S\ref{sec:whole-slice-contact} gives \st{F}
\[
\begin{array}{c|ccccccccc}
 \rho&0&0.05&0.1&0.2&0.3&0.5&0.7&0.9&1\\\hline
 I(1,0,1/4,\rho)&.93520&.93519&.93516&.93503&.93480&.93407&.93293&.93130&.93027\\
 I(1,.185,.174,\rho)&.83657&.83656&.83654&.83646&.83633&.83590&.83520&.83415&.83346\\
 \eta\text{-threshold for }500/701&.8941&.8939&.8934&.8912&.8876&.8756&.8666&.8666&.8666\\
 \text{same, rows re-optimised}&.8941&.8939&.8934&.8911&.8873&.8748&.8536&.8208&.8014
\end{array}
\]
Both integrals are even in $\rho$ ($X\mapsto-X$) and decrease
monotonically from $\rho=0$; so for this proof scheme the pure mixed
monomial is the worst point of the slice, and the concave majorant
\eqref{eq:slice-concave-relaxation} loses little
($0.71315$ certified against $0.7129$ float at $\eta=9/10$). At $\rho=1$
the re-optimised threshold $0.8014$ reproduces the $4/5$ of the diagonal
cone (\S\ref{sec:contact-concentration-cone}) from an independent
implementation. For $\rho\ge0.6$ the small-$u$ branch $L_\eta(t_+)$
binds at $\eta=9/10$ until $u_0$ is re-optimised.

\subsubsection{The reach of the support-cone method}
The same two rows applied to the candidate worst unit test of a sum of
$N$ orthogonal slices $\mathcal S_{X_1}+\dots+\mathcal S_{X_N}$, the
equal mixed sum $\psi_N=N^{-1/2}\sum_jX_j^2\bar Y_j$ (conditionally
Gaussian; a random search over unit elements of
$\mathcal S_{X_1}+\mathcal S_{X_2}$ found no element with a larger
support integral), give \st{F}
\[
\begin{array}{c|cccccc}
 N&1&2&3&4&8&\text{Gaussian}\\\hline
 I(1,.185,.174)&.83654&.83808&.83870&.83906&.83969&.84042\\
 \eta\text{-threshold, GPT rows}&.893&.945&.968&.981&\text{none}&\text{none}\\
 \eta\text{-threshold, rows re-optimised}&.893&.945&.968&.978&.996&\text{none}
\end{array}
\]
(``none'': this tested envelope is at least $500/701$ even at
$\eta=1$.) The certified two-slice theorem of \S\ref{sec:two-slice-contact} (shares $39/40$ and $49/50$) sits above this float threshold $0.945$, as its phase-averaged relaxation must. These are exploratory optimizations, not a theorem limiting
support certificates to four slices or requiring law-only bounds beyond
eight. Indeed, combining many support rows pointwise ($9\times10$
pairs $(b,d)$ and eight feasibility rows, the penalties' maximum taken
for the same response) certifies the fully diffuse Gaussian test at
share $\eta=1$: $\sup_{t,u}\mathcal N\le0.713172<500/701$ \st{F} (a $19\times21$ grid of rows; the family is near its limit),
and the mixed sums of $N=2,4,8$ slices down to $\eta=0.95,\,0.98,\,0.99$;
what remains structural is the loss at $\eta<1$, where the remainder
orthogonal to the single test is charged $p(1-\eta)V$ with no
penalty, so that a second test aligned with the remainder, or a law
row for it, is required; with the test taken to be the response's own
normalised cubic ($\eta=1$) the scheme reduces the central rows to
fixed-norm Gaussian domination of the PAIRED support
$\frac12\big(\E|Z(1+b(T-2))+d\psi|+\E|Z(1+b(T-2))-d\psi|\big)$ over unit
cubics $\psi$, visible sectors included (the one-sided form is false for visible-dependent $\psi$; see the correction below), for the finitely many $(b,d)$
rows; a uniform excess of about $5.5\cdot10^{-5}$ per row is tolerated (binding rows $(b,d)\approx(0.14\text{--}0.20,\,0.16\text{--}0.18)$ at $(t,u)\approx(0.845,0.284)$) \st{F}. Theorem~\ref{thm:sharp-decoupled-support} now gives a sharp
phase-preserving support for a whole orthogonal-input tensor block,
and Theorem~\ref{thm:phase-preserving-contact} reaches 24 holomorphic
inputs with two different certified support rows. That class includes
all mixed holomorphic quadratic terms, but requires the antiholomorphic
inputs to be orthogonal to them. The general repeated-slice space also
contains self-interaction and is a different class. The remaining task
is an exhaustive response cover, not an extrapolation from the sampled
thresholds in this table.

\subsubsection{An exact real-part reformulation of the remote functional}
\label{sec:real-part-reformulation}
The following identity was supplied on 22 October by a second
assistant answering the Conjecture-K prompt; the proof is complete.
\begin{lemma}[\st{P}: real-part form]\label{lem:real-part-form}
For every $G\in\Hc21(\C^n)$,
\[
 \Lambda(G)=2\E|L|+\pi\,\E\big(|\Re G|-|\Re z_1|\,a(|z_1|^2)\big)_+-\frac{\|G\|^2}{\epsilon},
 \qquad a(r)=\ell+\beta(r-2)>0 .
\]
Consequently Conjecture~K is equivalent to
$\pi\,\E(|\Re G|-|\Re L|)_+-\|G\|^2/\epsilon\le s_0:=S-2\E|L|
=0.00207424506\ldots$, and, writing
$(|x|-c)_+=\max_{|\tau|\le1}(\tau x-|\tau|c)$ and maximising over $G$
first, to
\[
 \sup_{G}\Lambda(G)-2\E|L|
 =\sup\Big\{\frac{\pi^2\epsilon}4\|\Pc21t\|_2^2-\pi\,\E\big(|t|\,|\Re L|\big):
 \ t:\C^n\to[-1,1]\text{ measurable}\Big\},
\]
where the supremum may be restricted to $t\in\{-1,0,1\}$-valued
functions.
\end{lemma}
\begin{proof}
Under the simultaneous rotation $z\mapsto e^{i\varphi}z$ both $L$
(charge $1$) and $G$ (bidegree $(2,1)$) pick up the factor
$e^{i\varphi}$, so $L\pm G$ are circularly symmetric and
$\E|L\pm G|=\frac\pi2\E|\Re(L\pm G)|$. For real $a,b$,
$|a+b|+|a-b|=2\max(|a|,|b|)=2|a|+2(|b|-|a|)_+$. Hence
$\E|L+G|+\E|L-G|=\pi\E|\Re L|+\pi\E(|\Re G|-|\Re L|)_+$ and
$\pi\E|\Re L|=2\E|L|$; finally $\Re L=\Re z_1\,a(|z_1|^2)$.
For the second form, $\E(\tau\Re G)=\Re\langle\tau,G\rangle
\le\|\Pc21\tau\|\,\|G\|$ with equality at
$G\propto\Pc21\tau$, and $\sup_{g\ge0}(\pi\|\Pc21\tau\|g-g^2/\epsilon)
=\pi^2\epsilon\|\Pc21\tau\|^2/4$; for fixed $G$ the maximising $\tau$ is
$\operatorname{sign}(\Re G)\,\mathbf 1\{|\Re G|>|\Re L|\}$.
\end{proof}
Numerically, at the Gaussian point the gain term is
$\pi\E(|\Re G|-|\Re L|)_+=0.05541$ and the identity was reproduced to
$1.4\cdot10^{-6}$ by direct quadrature \st{F}; the whole gain is
produced near the hyperplane $\Re z_1=0$. The second form shows how
sharp any proof must be: $1.3\cdot10^{-4}$ is $0.23\%$ of the gain
term. It also recasts Conjecture~K as a statement about one bounded real
function, its cubic projection energy against its $|\Re L|$-weighted
mass, which is the shape of the original-response programme of this
chapter (bounded responses, cubic projections and retained costs).

Two further points from the same source, recorded with their status.
(i) \st{F} In its local analysis the worst gain/cost ratio at the
Gaussian point is $0.852$, in the direction $iz_1Y$ with $Y$ real in
$\Hc11(w)$; this is the imaginary Y-sector of the table above
($1-0.395/2.644=0.85$), an independent confirmation. Its global searches
over $\Hc21(\C^2)$ and $\Hc21(\C^3)$ plus a Gaussian part returned to
$\sqrt{v^*}U$ within quadrature error. (ii) \st{F} A semidefinite
computation of the best HJM-type profile-matrix bound at the extremal
configuration loses at least $2.19\cdot10^{-4}$ for $\Lambda$ as stated
(more than the $1.3\cdot10^{-4}$ budget) but only $8.8\cdot10^{-5}$ once
the retained term $-\|\Pc10k\|^2$ of the response functional is kept,
leaving about $4\cdot10^{-5}$. This agrees with the diagnosis of
\S\ref{sec:gaussian-point-cone-reach} and of the certificate audit
(the row/gate structure, not the law relaxation, carries most of the
loss), and it says again that the target is the response functional
$\mathcal N$ with its negative squares, not raw $\Lambda$.

\paragraph{Caution.} The same source proposes to settle the
$z_1$-independent case from the concavity of $\Phi_0(s)=\E|L+\sqrt s|$
(proved in \S\ref{sec:whole-slice-contact}) together with the convex
order $|q|^2\ge_{\rm cx}\|q\|^2\,\mathrm{Exp}(1)$ on $\Hc21$. The
first fact is true and the implication is correct, but the convex order
is \emph{false}: see the negative register (stop-loss conjecture CX,
$q_\theta=z_0[\cos\theta(|z_0|^2-2)/\sqrt2+i\sin\theta(|z_1|^2-1)]$,
density of $|q|^2$ at $0$ equal to $0.85<1$) and the certified Laplace
counterexample $\E e^{-30|H|^2}<1/31$ in Chapter~\ref{ch:K-state}. The
$z_1$-independent case therefore has to go through the true constraints
on the law of $|q|^2$, the moment bounds and the Stein cuts, which is
exactly what the Lean rows do.


\subsubsection{Reported: a uniform local theorem and an improper Gaussian-mixture theorem}
\label{sec:local-stability-reported}
The second assistant subsequently proved (22 October) a local theorem
in two forms. The first form (package
\texttt{cubic-chaos-local-stability.zip}, not received by the IDE;
\st{P} as reported) is the following. With
$K=G-q$, $Q=\|q\|^2$, $D=\|K\|^2$ and
$\mathcal J_s(G)=\E|L+\sqrt sU+G|+\E|L+\sqrt sU-G|-p(s+\|G\|^2)$
($U$ a fresh standard complex Gaussian):
\[
 \tfrac12\tfrac{d^2}{dt^2}\mathcal J_s(tG)\big|_{t=0}\le F'(s)Q-\tfrac{D}{40}
 \quad(s>0),\qquad
 \|G\|\le10^{-6}\ \Longrightarrow\ \mathcal J_{v^*}(G)\le S-\tfrac{D}{80}-\tfrac{Q^2}4 .
\]
So at $s=v^*$ the kernel of the Hessian is exactly the $Z$-independent
cubics, in every dimension and with arbitrary coordinate overlap among
$q,u,v,A,B,c$. The A--B coupling coefficient is
$d=-2^{-1/2}\int_0^\infty e^{-t}t(t-1)\,d_s(R(t))\,dt$ with
$d_s=f_s'/r-f_s''$, $0\le d_s\le1/r$, bounded by $|d|<23/50$ (split at
$t=1$), and the block
$\begin{pmatrix}\lambda-123/100&-23/100\\-23/100&\lambda-12/25\end{pmatrix}$,
$\lambda=p-1/40$, has positive determinant $53919007/28607048000$. The
fourth-order term along $q$ is $\frac14M''(v^*)\E|q|^4<-\frac12Q^2$
($M''(v^*)<-2$ by interval arithmetic; $\E|q|^4\ge2Q^2$), and the
remainder uses the dimension-free Wick bounds $\E|J|^4\le180\|J\|^4$,
$\E|J|^6\le45360\|J\|^6$. A genuine-cubic consequence: with the
disjoint-product block $B_{m,v}$,
$|\Lambda(G+B_{m,v})-\mathcal J_v(G)|\le3\sqrt{\pi v}/(2m)$, hence
$\Lambda(G+B_{m,v^*})<S$ for $0<\|G\|\le10^{-6}$ and
$m\ge\lceil3/(5\Delta)\rceil$, $\Delta=D/80+Q^2/4$.
Separately, an \emph{improper} Gaussian-mixture theorem: if $q$ is
independent of $Z$ and conditionally (on a $\sigma$-field independent
of $Z$) a centred real two-dimensional Gaussian with arbitrary
covariance $(\lambda_1,\lambda_2)$, then $\Lambda(q)\le S$; proof:
$h(t)=2\E|L+\sqrt t|$ is concave (radially decreasing density of $L$
and the radial Poisson equation), and exchanging two independent real
Gaussians gives
$\E h(\lambda_1N_1^2+\lambda_2N_2^2)\le\E h(\tfrac{\lambda_1+\lambda_2}2(N_1^2+N_2^2))=M(\lambda_1+\lambda_2)$.
This covers $q=\sum_j(Y_jA_j(X)+\bar Y_jB_j(X))$ with the same $Y_j$
conjugated and unconjugated and overlapping quadratic coefficients,
extending the proper case of the cubic-chaos package. The IDE checked
the arguments as written (the exchange step and the frequency
computation above are elementary).

The second form is interval-certified and was replayed by the IDE
\st{P,C} (\texttt{verify\_local\_stability.py}, Arb ball arithmetic at
100 bits, archived with its logs in
\texttt{conjecture-K-report/gpt-packages/local-stability-20261022/}).
With $m_s(x)=\E|x+\sqrt sU|$, whose Hessian has radial and tangential
eigenvalues $R_s=C_se^{-x}(I_0-I_1)(x)$, $D_s=C_se^{-x}(I_0+I_1)(x)$,
$C_s=\sqrt\pi/(2\sqrt s)$, $x=r^2/(2s)$, $h_\pm=(R_s\pm D_s)/2$, the
$t^2$-coefficient of $\E m_s(L+tH)+\E m_s(L-tH)-p(s+t^2M)$ is
$Q_s(H)=\E[h_+(|L|)|H|^2+h_-(|L|)\Re(e^{-2i\arg Z}H^2)]-pM$, and at
$s_0=403448/10^7$ the sector coefficients are enclosed as
\[
\begin{array}{c|ccccccc}
 &a_u&a_v&b_{uv}&a_A&a_B&a_{c,R}&a_{c,D}\\\hline
 \text{certified}&-0.46563&-0.93337&-0.15499&-0.19554&-0.70055&-1.17043&-0.42495\\
 \text{IDE float }h/2&-0.466&-0.933&\text{(not computed)}&-0.198&-0.700&-1.17&\text{(not computed)}
\end{array}
\]
($a_u=\E[(T-1)^2h_+]-p$, $a_v=\frac12\E[T^2h_+]-p$,
$b_{uv}=2^{-1/2}\E[T(T-1)h_-]$, $a_A=\E[TD_s]-p$, $a_B=\E[Th_+]-p$,
$a_{c,R}=\frac12\E[T(T-2)^2R_s]-p$, $a_{c,D}=\frac12\E[T(T-2)^2D_s]-p$;
the A-row uses $|\E A^2|\le\E|A|^2$, i.e.\ the worst, imaginary, Y-sector).
The IDE's float table is thus certified entry by entry. Consequences
proved there: $Q_{v^*}(H)\le-\frac{19}{100}d$ with
$d=\|H-\E[H\mid w]\|^2$ (transfer from $s_0$ to $v^*$ through
$\|\partial_sD^2m_s\|\le6s^{-3/2}$); with $\Psi_H(t)$ the value at total
variance $v^*$ split as $v^*-Mt^2$ Gaussian plus $tH$,
$Mt^2\le d/(200000M)\Rightarrow\Psi_H(t)\le S-\frac7{200}t^2d$, using
$\E|H|^4\le62M^2$ and $\|D^4_\xi m_s\|\le20s^{-3/2}$; and for
$Z$-independent $q$ the quartic coefficient
$\frac14F''(v^*)(\E|q|^4-2M^2)<0$ strictly ($F''<0$ by the heat
equation and the radially decreasing density of $L$; $\E|q|^4>2M^2$ by
the OU product rule). It also certified a Laplace obstruction of the
same kind as Chapter~\ref{ch:K-state}'s:
$Q=X(4(|X|^2-2)+3i(|Y|^2-1))/\sqrt{41}$ has
$\E e^{-20|Q|^2}<0.0453<1/21=\E e^{-20|U|^2}$. These results are consistent with the sector table: they
turn its float margins into a rigorous, if conservative ($1/40$),
uniform gap, and they confirm that the only flat direction at the
benchmark is the Gaussian family.
\subsubsection{Concentrated remote laws are far from the benchmark}
Two float scans \st{F} on the droplet complete the picture. Along the
CLT family $G_N=(2N)^{-1/2}\sum_{j\le N}X_j(|X_j|^2-2)$ the optimum
$\sup_s\Lambda(sG_N)$ increases monotonically to $S$:
$S-1.95\cdot10^{-3}$, $-1.85$, $-1.74$, $-1.67$, $-1.53$,
$-1.30\cdot10^{-3}$ for $N=1,2,3,4,6,10$, with preferred remote energy
$0.0026,\,0.0048,\,0.0069,\,0.0086,\,0.0114,\,0.0159$ against
$v^*=0.0403$. Along the slice $\psi_\rho$ the remote-only optimum is
$S-1.68\cdot10^{-3}$ at $\rho=0$ (the conditionally Gaussian end),
decreasing to $S-2.02\cdot10^{-3}$ at $\rho=1$, with preferred energy
$0.005$--$0.008$. Adding a coupled A-sector $\alpha(T-1)X$ on the slice's own driver $X$ and optimising $(s,\alpha)$ gains at most $8\cdot10^{-5}$ (at $\rho=1$, $\alpha\approx0.01$ real; $1.5\cdot10^{-5}$ at $\rho=0.5$, nothing at $\rho=0$; every $|\alpha|\ge0.03$ and every imaginary $\alpha$ loses), i.e.\ at most $4\%$ of the margin. So concentrated cubics gain about $3\cdot10^{-4}$
over $2\E|L|$ where the Gaussian gains $2.07\cdot10^{-3}$: the
conjectured maximiser is approached only through diffuseness, slowly
($N=10$ is still $1.3\cdot10^{-3}$ short), and the construction bound
$F(v)-6v^2/m\le\Lambda(G_{m,v})$ of the cubic-chaos package is loose by
a factor of about five at small $m$. The equal mixed sums $\psi_N=N^{-1/2}\sum_{j\le N}X_j^2\bar Y_j$ (conditionally Gaussian) approach $S$ slowly:
$S-\sup_s\Lambda(s\psi_N)=1.68,\,1.42,\,1.22,\,1.08,\,0.90,\,0.73,\,0.55,\,0.43,\,0.32\cdot10^{-3}$
for $N=1,2,3,4,6,8,12,16,24$, with $v^*$ rising from $0.008$ to $0.034$. For the response programme this
means that the concentrated branch of the cover has room of order
$10^{-3}$ in $\Lambda$, that the band of five to eight slices where the two-row support method fails still has
$7$--$9\cdot10^{-4}$ of room (five to seven times the whole $1.3\cdot10^{-4}$ allowance, so it needs a certificate that is merely not
too lossy, not a sharp one), and that only genuinely diffuse laws leave no room at all.

\subsubsection{Moment-only bounds for the independent sector have a floor}
\label{sec:independent-sector-floor}
The second assistant also proved \st{P} that $h(x)=2\E|L+\sqrt x|$ is a
Bernstein function (a positive mixture of $1-e^{-\lambda x}$), so the
three-moment relaxation $\max\E h(X)$ subject to $\E X=v$,
$\E X^2=mv^2$, $\E X^3\le nv^3$ is attained by an explicit two-point
law with atoms $vr_\pm$, $r_\pm=\frac12\big(A\pm\sqrt{(A-2)^2+4(m-1)}\big)$,
$A=(n-m)/(m-1)$; with the sixth-moment coefficient $1977/40$ it certified
$\Lambda(q)<S+1.05\cdot10^{-5}$ for $q$ independent of $z_1$ \st{P,C}.
Evaluating the same two-point formula with the sharper circular constant
$C_6=40.45313$ (all of whose tensor facts are now proved in Lean) moves
the worst case to the Gaussian-kurtosis boundary $m\to2$ and gives
$S+9.76\cdot10^{-6}$ \st{F}: only $3\%$ lower, because at $m=2$ a two-point
law with the first three moments of $\mathrm{Exp}(1)$ already beats the
exponential law by about $10^{-5}$, independently of $C_6$. Moment-only
relaxations of the independent sector therefore stall about $10^{-5}$
above $S$; below that, chaos-specific constraints near $m=2$ are needed,
which is what the Stein, spectral and contraction-Gram cuts of the
complete certificate supply. This has now been done: \S\ref{sec:independent-stein-36e-8} certifies $\Lambda(q)<S+3.6\cdot10^{-7}$ \st{P,C} for every independent cubic, from the sixth-moment theorem, the radial Stein identity, the separate spectral budgets and the covariance order (IDE replay PASS).

An independent audit (22 October, package
\texttt{cubic-chaos-independent-bound-audited.zip}, replayed) adds three
facts \st{P,C}. (i) The independent-sector bound survives an independent
Gaussian summand: $\Lambda(q+gU)<343491/250000<S+1.05\cdot10^{-5}$ for every
independent $q$ and $g\ge0$, because with $r=A/(A+B)$ ($A=\E|q|^2$,
$B=g^2$) the normalised moments satisfy $m_W-2=r^2(m_q-2)$ and
$n_W-9m_W+12=r^3(n_q-9m_q+12)$, so the sixth-moment constant is unchanged.
(ii) With $r=\|G-\E[G\mid w]\|$, $\Lambda(G)<1.373964+2r-pr^2$; hence the
relaxed bound $S+1.3\cdot10^{-4}$ holds on the tube $r\le1.19\cdot10^{-4}$
around the whole independent subspace. (iii) The seven second-variation
coefficients of \S\ref{sec:local-stability-reported} hold uniformly for
$v\in[0.040344,0.040346]$ (a derivative bound
$|\partial_sK_\nu|\le\sqrt\pi/(4s^{3/2})$ with allowance $2.256\cdot10^{-4}$),
so they are certified at $v^*$ itself.


\paragraph{Correction: the general support test must retain both signs
(22 October, \st{P,C}).}
The IDE's 07:40 multi-row observation cannot use global negation to
replace paired support by one modulus when the test depends on $Z$.
Global negation flips both $L$ and the cubic and preserves each modulus
separately. The one-sided Gaussian domination it proposed is false:
with the peak $L=Z(a+bT)$, $a=699/1250$, $b=179/1250$,
$H=Z(T-2)/\sqrt2$ and $d=1/100$, positivity of the two radial
coefficients gives exactly
\[
 \E|L-dH|=\E|L|+\frac{d\sqrt\pi}{4\sqrt2},\qquad
 \frac{\E|L-dH|+\E|L+dH|}{2}=\E|L|.
\]
For a fresh proper Gaussian $W$, concavity in its variance gives
\[
 \E|L+dW|\le\E|L|+\frac{d^2}{4}\E|L|^{-1}
             \le\E|L|+\frac{d^2\sqrt\pi}{4a}.
\]
Thus the one-sided gap exceeds
$(17/40)((7/10)d-d^2/a)=2533/873750>1/400$.
The rational checker in
\pathref{spherical-response-contact-20261022/} verifies the constants.
This is not a counterexample to K or to the paired support inequality.
The latter remains the correct target for visible-dependent tests.
All contact theorems above reduce to one side only for remote tests,
where negating remote coordinates leaves $L$ fixed.
